\itemhead{3.2 思考 · 未编号题组（教材p.83）}{XJ-B1-C03-TH-P083-Y0867}
\label{XJ-B1-C03-TH-P083-Y0867}
{\footnotesize 湘教版教材p.83 / PDF第90页 / 3.2 思考 / 未编号题组}
\field{原题}{如何用数学符号语言来描述函数的奇偶性？}
\field{参考答案}{先要求定义域 $D$ 关于原点对称；若 $\forall x\in D,f(-x)=f(x)$，则为偶函数；若 $f(-x)=-f(x)$，则为奇函数。}
{\small \field{关联知识}{F08 奇偶性、对称与函数性质应用（教材帮p.113、114、115 / PDF122、123、124）}
\field{方法与易错点}{检查定义域，再代入 $-x$；求另一侧解析式时同时写出 $-x$ 所在区间。 易错：只看公式不看定义域；用一对点证明奇偶；负半轴代换遗漏最外层负号。}
\field{答案依据}{答案册未提供本栏目的独立对应答案；依据原题、教材概念补写参考回应。开放性总结不设唯一数值答案。}}
